\chapter{Introduction}
\label{ch:introduction}

\section{Le risque climatique est devenu une ligne de compte}

Un assureur dommages français ne peut plus tarifer un bien immobilier sans
regarder ce qui lui arrive : submersion, ruissellement, mouvements de terrain,
retrait-gonflement des argiles, séisme, feux de forêt. Un promoteur ne peut plus
livrer un programme sans répondre à la même question, posée par son assureur,
son banquier ou son acquéreur.

Les données existent. elles sont publiques, souvent de grande qualité, et
produites par des acteurs identifiés : Géorisques (BRGM, MTE), la BDNB (CSTB),
le géocodage de la Géoplateforme (IGN), l'hydrographie de la BD TOPO, Open-Meteo,
Panoramax. Le problème n'est donc pas l'absence de données : c'est leur
\textbf{dispersion}, leur \textbf{hétérogénéité} et l'absence de clé de jointure
commune.

\begin{encadre}[title={Le décalage central}]
Les référentiels publics décrivent le risque à l'échelle de la \textbf{commune}
ou du \textbf{secteur}. Une décision d'assurance ou d'investissement se prend à
l'échelle du \textbf{bâtiment}. Tout l'écart entre les deux est du travail manuel
— et, quand il est franchi sans méthode, une source d'erreur silencieuse.
\end{encadre}

\section{Objet du document}

Ce rapport décrit \textbf{Typhoon}, un service dont la fonction est de combler
cet écart : à partir d'une adresse, produire un enregistrement canonique qui
joint les aléas du référentiel Géorisques au bâtiment lui-même, en déclarant
pour chaque fait son origine et sa finesse.

Le document est organisé en trois mouvements, dans l'ordre du problème vers la
solution :

\begin{itemize}
  \item \textbf{le problème métier} (\S\ref{ch:douleurs}) — les douleurs que
    subissent concrètement l'assureur et le promoteur, formulées en coûts et en
    risques, et non en fonctionnalités ;
  \item \textbf{la réponse} (\S\ref{ch:reponse} à \S\ref{ch:robustesse}) — le
    contrat de données, l'architecture en diagrammes UML, les parcours réels et
    les garanties de robustesse ;
  \item \textbf{le bilan} (\S\ref{ch:resultats} et \S\ref{ch:douleurs-restantes})
    — ce qui a été mesuré, puis les points de douleur qui \emph{subsistent},
    chacun suivi de son remède.
\end{itemize}

\section{Les acteurs}

Quatre profils utilisent le service, avec des besoins différents qui expliquent
la forme du contrat de données.

\begin{itemize}
  \item \textbf{L'assureur} — tarifie un portefeuille, doit justifier sa
    décision devant son régulateur et ne veut pas d'une boîte noire. Il lui faut
    la provenance de chaque fait et le droit de ne pas conclure quand la source
    n'a pas répondu.
  \item \textbf{Le promoteur} — engage un programme immobilier et doit répondre
    vite à une question de risque, en amont de l'acquisition.
  \item \textbf{Le gestionnaire de portefeuille} — suit un ensemble de biens
    dans le temps et compare ce qui est comparable.
  \item \textbf{L'analyste} — ouvre la fiche bâtiment, lit les champs BDNB
    verbatim et cherche ce qui motive un écart de tarification.
\end{itemize}

\section{Périmètre et non-périmètre}

\begin{tabularx}{\textwidth}{@{}lX@{}}
\toprule
\textbf{Dans le périmètre} & \textbf{Hors périmètre (par décision)} \\
\midrule
Diagnostic par bâtiment, France métropolitaine & Calcul de prime, tarification actuarielle \\
Aléas Géorisques + classe TRI au point & Score composite, bande de risque, notation \\
Vulnérabilité BDNB verbatim (139 champs) & Recommandation générée présentée comme un fait \\
Provenance, résolution, attribution sur chaque fait & Prédiction de sinistre \\
Estimation de dommages paramétrique (scénarios TRI) & Données Copernicus dans le contrat sérialisé \\
\bottomrule
\end{tabularx}

\begin{invariant}
\textbf{Aucun score, aucune bande, aucune recommandation générée} n'entre dans
le contrat sérialisé. Ce n'est pas une réserve commerciale, c'est une contrainte
technique : le schéma canonique du backend
(\code{diagnostic\_record.py}) \emph{refuse} ces champs à la validation. Une
donnée qui n'a pas de source ne peut pas être ajoutée au contrat sans qu'un
amendement explicite soit voté.
\end{invariant}
